
and let \(x\) be a point of \(F_2\).

\begin{enumerate}
\item[\SGAInlineRefTarget{lem:sga3-V-7.4:item:i}(i)] If \(u\) is flat, \(d'\) is flat at \(x\) if and only if
  \(d\) is flat at \(v(x)\).
\item[(ii)] If \(d\) is locally of finite type, \(d'\) is quasi-finite
  at \(x\) if and only if \(d\) is quasi-finite at \(v(x)\).%
  \footnote{\label{note:sga3-V-42}Editor's note: The conditions are
  sufficient by base change (EGA~II, 6.2.4(iii), and
  EGA~IV\(_2\), 2.1.4). Conversely, put
  \(y=d'(x)\) and \(z=u(y)=d(v(x))\), and suppose that \(d'\) is flat
  at \(x\) and that \(u\) (hence also \(v\)) is flat. Then
  \(\mathcal O_{v(x)}\to\mathcal O_x\) is faithfully flat, as is
  \(\mathcal O_z\to\mathcal O_y\to\mathcal O_x\). Consequently,
  \(\mathcal O_z\to\mathcal O_{v(x)}\) is faithfully flat
  (EGA~IV\(_2\), 2.2.11(iv)). Finally, suppose that \(d\) is locally
  of finite type and that \(d'\) is quasi-finite at \(x\). Then
  \(v(x)\) is isolated in its fiber \(d^{-1}(z)\), since \(x\) is
  isolated in its fiber
  \(d'^{-1}(y)=d^{-1}(z)\otimes_{\kappa(z)}\kappa(y)\). By Chevalley's
  semicontinuity theorem there is therefore an open neighborhood of
  \(v(x)\) every point of which is isolated in its fiber
  (EGA~IV\(_3\), 13.1.3 and 13.1.4), and hence \(d\) is quasi-finite
  at \(v(x)\).}
\end{enumerate}

We have thus proved that \(X_0\) has a covering by saturated open
subsets \(W\) such that the groupoid \(W_\ast\) induced by \(X_\ast\)
on \(W\) possesses a quasi-section.%
\footnote{\label{note:sga3-V-43}Editor's note: The sequel has been
modified to take advantage of the additions made to
Lemma~\hyperref[lem:sga3-V-6.1]{6.1}.}

By Lemma~\hyperref[lem:sga3-V-6.1]{6.1} and the reductions stated
after Theorem~\hyperref[thm:sga3-V-7.1]{7.1}, this proves assertions
\hyperref[thm:sga3-V-7.1-i]{(i)} and
\hyperref[thm:sga3-V-7.1-iii]{(iii)} of that theorem, as well as the
facts that \(p\) is surjective and open and that \(p\) and \(Y\to S\)
are locally of finite presentation. Moreover, because \(X_0\to S\)
is quasi-projective, and hence separated and of finite type, \(p\) is
separated; the proof of Lemma~\hyperref[lem:sga3-V-6.1-ii]{6.1(ii)}
shows that \(p\) and \(Y\to S\) are of finite presentation.

To prove that \(p\) is proper, it remains to prove that it is
universally closed. Since this assertion is local on \(Y\), we may
work over a saturated open subset \(W\) such that the groupoid
\(W_\ast\) induced by \(X_\ast\) on \(W\) possesses a quasi-section
\(U\) (because \(X_0\) is covered by such open subsets). With the
notation of \hyperref[par:sga3-V-6-a]{6(a)}, there is a commutative
diagram
\SGAThreeDiagram[0.48\textwidth]
  {figures/exp5/exp5_localp28_WVU_Z_commutative.png}
  {The properness diagram over the open subset \(Z\subseteq Y\).}
  {fig:sga3-V-7.1-properness}
  {LNM 151, p.~281}
where \(Z\) is an open subset of \(Y\), all the arrows are
surjective, and \(q\) is integral. By hypothesis
\(d_0:X_1\to X_0\) is proper, so \(u\), which is obtained from \(d_0\)
by base change, is proper as well. Consequently \(r\) is universally
closed, and hence so is \(p\), since \(v\) is surjective.

Finally, because \(p\) is surjective and universally closed and
\(X_0\) is quasi-projective, hence separated, the diagonal
\(\Delta_{Y/S}(Y)\) is closed in \(Y\times_S Y\): it is the image under
\(p\times p\) of the diagonal \(\Delta_{X_0/S}(X_0)\). Thus \(Y\) is
separated over \(S\). This completes the proof of
Theorem~\hyperref[thm:sga3-V-7.1-ii]{7.1(ii)}.

The assertions in Theorem~\hyperref[thm:sga3-V-7.1-iv]{7.1(iv)} are
local on \(Y\). Since the saturated open subsets \(U_z\) cover
\(X_0\), it is enough to verify them after replacing \(X\) and \(Y\)
by \(U_z\) and \(V_z=p(U_z)\), respectively. As observed at the
beginning of the proof of Theorem~\hyperref[thm:sga3-V-7.1]{7.1},
Lemmas~\hyperref[lem:sga3-V-1.1]{1.1},
\hyperref[sga3:exp5:lemma:1.2]{1.2}, and
\hyperref[lem:sga3-V-6.1-i]{6.1(i)} show that
\((V_z,p|_{U_z})\) is the cokernel, both in \((\mathbf{Sch})\) and in
the category of ringed spaces, of the groupoid induced by \(X_\ast\)
on \(U_z\). The hypothesis that \(d=(d_0,d_1)\) is a monomorphism is
preserved under the base change \(U_z\to X_0\). The first two
assertions of Theorem~\hyperref[thm:sga3-V-7.1-iv]{7.1(iv)} therefore
follow from Lemma~\hyperref[lem:sga3-V-6.1-iv]{6.1(iv)}.

Let us finally prove the consequences stated at the end of \hyperref[thm:sga3-V-7.1-iv]{part~(iv)}
(see \(\hyperref[bib:sga3-V-Ray67a]{[\mathrm{Ray67a}]}\), Theorem~1(iii)). By hypothesis, the
groupoid \(X_\ast\) comes from an equivalence relation
\(R\to X_0\times X_0\). We have proved that \(R\) is effective
(\SGARefLink{sga3:iv:definition:3:3:2}{Expos\'e~IV, 3.3.2}) and that
\(p:X_0\to Y=X_0/R\) is faithfully flat and of finite presentation.
Consequently, if \((M)\) denotes the family of faithfully flat
morphisms locally of finite presentation, then \(R\) is
\((M)\)-effective. By
\SGARefLink{sga3:iv:corollary:6:3:3}{Expos\'e~IV, 6.3.3}, the pair \((Y,p)\) therefore
represents the quotient sheaf of \(X_0\) by \(R\) for the fppf
topology, and the assertions concerning base change follow from
\SGARefLink{sga3:iv:proposition:3:4:3:1}{Expos\'e~IV, 3.4.3.1}.

\section*{8. Passage to the quotient by a flat groupoid that is not
  necessarily proper}
\phantomsection\label{sec:sga3-V-8}
\addcontentsline{toc}{section}{8. Passage to the quotient by a flat groupoid that is not necessarily proper}

\medskip
\phantomsection\label{thm:sga3-V-8.1}
\noindent\textbf{Theorem 8.1.}---%
\footnote{\label{note:sga3-V-44}Editor's note: There is a largest open
subset \(W\) of \(X_0\) satisfying the conclusions of \hyperref[thm:sga3-V-8.1]{the theorem}.
Indeed, let \(W\) be an open subset as in \hyperref[thm:sga3-V-8.1]{the theorem} and let
\(W^\sharp\) be a dense saturated open subset of \(W\). Since \(p\)
is open, \(V^\sharp=p(W^\sharp)\) is open in \(V\), and
\(W^\sharp=p^{-1}(V^\sharp)\), because \(W^\sharp\) is saturated. By
Lemma~\hyperref[lem:sga3-V-1.1]{1.1}, \(V^\sharp\) is a cokernel for
the groupoid induced on \(W^\sharp\). Thus two open subsets \(W\) and
\(W'\) satisfying the conclusions of \hyperref[thm:sga3-V-8.1]{the theorem} may be glued along
their intersection \(W^\sharp\), and conditions \hyperref[thm:sga3-V-8.1-i]{(i)}--\hyperref[thm:sga3-V-8.1-iv]{(iv)}, together
with the property that \(p\) and \(V\to S\) are locally of finite
presentation, are preserved. Conclusion (ii') follows, as in the
proof of Lemma~\hyperref[lem:sga3-V-6.1-ii]{6.1(ii)}, from the
hypothesis that \(X_0\) is of finite type over the Noetherian scheme
\(S\).

Lemmas~\hyperref[lem:sga3-V-1.1]{1.1} and
\hyperref[sga3:exp5:lemma:1.2]{1.2} also show that the union of all
saturated open subsets \(U\subseteq X_0\) such that the open subset
\(p(U)\subseteq Y\) is a scheme and
\(p|_U:U\to p(U)\) is a morphism of schemes is the largest saturated
open subset \(\Omega\subseteq X_0\) satisfying \hyperref[thm:sga3-V-8.1-i]{condition~(i)} of
Theorem~\hyperref[thm:sga3-V-8.1]{8.1}. \hyperref[thm:sga3-V-8.1]{The theorem} shows that
\(\Omega\) contains a dense open subset \(W\), but it is not immediate
that \(\Omega\) satisfies \hyperref[thm:sga3-V-8.1-ii]{properties (ii)}--\hyperref[thm:sga3-V-8.1-iv]{(iv)}. For more precise
results, and for the representability of the fppf quotient sheaf
\(\widetilde{X/R}\) (where \(R\) denotes the groupoid \(X_\ast\) with
base \(X=X_0\)), see \(\hyperref[bib:sga3-V-Ray67a]{[\mathrm{Ray67a}]}\), \(\hyperref[bib:sga3-V-Ray67b]{[\mathrm{Ray67b}]}\), and
Appendix~I of \(\hyperref[bib:sga3-V-An73]{[\mathrm{An73}]}\). These works use weaker hypotheses:
\(S\) may be arbitrary, \(X\) is locally of finite type over \(S\),
and \(R\) is an \(S\)-groupoid with base \(X\) such that \(d_0\) (and
hence \(d_1\)) is flat and of finite presentation. In particular, if
\(\widetilde{X/R}\) is represented by an \(S\)-scheme \(Y\), then \(Y\)
is also the cokernel in the category of ringed spaces. The converse
is false in general (see \(\hyperref[bib:sga3-V-Mum65]{[\mathrm{Mum65}]}\), Chapter~0, \S3,
Example~0.4, cited in \(\hyperref[bib:sga3-V-Ray67a]{[\mathrm{Ray67a}]}\), Remark~1), but it is true
when \(d=(d_0,d_1)\) is an immersion. Under this hypothesis,
\(p:\Omega\to Z:=\Omega/R_\Omega\) is faithfully flat and of finite
presentation. If, in addition, \(S\) is locally Noetherian, then a
codimension-one point \(x\) of \(X\) belongs to \(\Omega\) if and only
if the graph of the groupoid induced on
\(\operatorname{Spec}(\mathcal O_{X,x})\) is closed. See
\(\hyperref[bib:sga3-V-Ray67a]{[\mathrm{Ray67a}]}\), Proposition~1; \(\hyperref[bib:sga3-V-Ray67b]{[\mathrm{Ray67b}]}\),
Proposition~1 and Theorems~2, 1, and 4; and \(\hyperref[bib:sga3-V-An73]{[\mathrm{An73}]}\),
Theorems~5 and 6 on pages~66--67 and Proposition~3.3.1 on page~49.
For an action of an algebraic group over an algebraically closed field
\(k\), see also \(\hyperref[bib:sga3-V-DR81]{[\mathrm{DR81}]}\).}%
\textit{Let \(S\) be a Noetherian scheme and}
\SGAThreeDiagram[0.57\textwidth]
  {figures/exp5/exp5_localp29_theorem81_groupoid.png}
  {The \((\mathbf{Sch}/S)\)-groupoid of \hyperref[thm:sga3-V-8.1]{Theorem~8.1}.}
  {fig:sga3-V-8.1-groupoid}
  {Polo--Gille re-edition, Expos\'e V}
\textit{a \((\mathbf{Sch}/S)\)-groupoid such that \(d_1\) is flat and
of finite type, \(X_0\) is of finite type over \(S\), and the morphism
\(X_1\to X_0\times_S X_0\) with components \(d_0\) and \(d_1\) is
quasi-finite.}

\textit{There exists a dense saturated open subset \(W\subseteq X_0\)
having the following properties:}
\begin{enumerate}
\item[(i)]\phantomsection\label{thm:sga3-V-8.1-i}%
  \textit{if
  \(W_2\overset{w'_i}{\rightrightarrows}
    W_1\overset{w_j}{\rightrightarrows}W\)
  is the groupoid induced by \(X_\ast\) on \(W\), then
  \((w_0,w_1)\) has a cokernel \((V,p)\) in
  \((\mathbf{Sch}/S)\); moreover, \((V,p)\) is a cokernel of
  \((w_0,w_1)\) in the category of all ringed spaces;}
\item[(ii)]\phantomsection\label{thm:sga3-V-8.1-ii}%
  \textit{\(p\) is surjective and open;}
\item[(ii')]\phantomsection\label{thm:sga3-V-8.1-iip}%
  \textit{\(p\) and \(V\to S\) are of finite presentation;}
\item[(iii)]\phantomsection\label{thm:sga3-V-8.1-iii}%
  \textit{the morphism \(W_1\to W\times_VW\) with components \(w_0\)
  and \(w_1\) is surjective;}
\item[(iv)]\phantomsection\label{thm:sga3-V-8.1-iv}%
  \textit{if \((d_0,d_1)\) is an equivalence pair, then
  \(W_1\to W\times_VW\) is an isomorphism and \(p\) is faithfully
  flat.}
\end{enumerate}

We shall prove that \(W\) may be chosen so that the
\((\mathbf{Sch}/S)\)-groupoid \(W_\ast\) induced by \(X_\ast\)
possesses a quasi-section (see \hyperref[sec:sga3-V-7]{\S7}).
Theorem~\hyperref[thm:sga3-V-8.1]{8.1} will then follow from
Lemma~\hyperref[lem:sga3-V-6.1]{6.1}.

Suppose provisionally that for every point \(z\in X_0\) closed
relative to \(S\) (see \hyperref[sec:sga3-V-7]{\S7}) there is a
saturated open subset \(W_z\) that possesses a quasi-section and meets
every irreducible component of \(X_0\) passing through \(z\). Then the
complement \(X_0-W_z\) is saturated: its saturation
\(d_1(d_0^{-1}(X_0-W_z))\) is open and does not meet \(W_z\). If this
complement is nonempty, choose in it a point \(z'\) closed relative to
\(S\), and associate to \(z'\) an open subset \(W_{z'}\) as above. We
may further suppose that \(W_{z'}\subseteq X_0-W_z\). Then \(W_z\) and
\(W_{z'}\) are disjoint, and the groupoid induced by \(X_\ast\) on
\(W_z\cup W_{z'}\) possesses a quasi-section. The process must stop,
because \(X_0\) has only finitely many irreducible components. It
therefore remains to construct \(W_z\).

For this purpose we may suppose that \(S\) is affine. Let \(y\in X_1\)
satisfy \(d_1(y)=z\), let \(X\) be an affine open subset of \(X_0\)
containing \(d_0(y)\), let \(Y\) be the inverse image of \(X\) under
\(d_0:X_1\to X_0\), and let \(u:Y\to X\) and \(v:Y\to X_0\) be the
morphisms induced by \(d_0\) and \(d_1\). Since \(X\) is affine, hence
quasi-projective, Lemma~\hyperref[lem:sga3-V-7.2]{7.2} provides a
subscheme \(F\subseteq X_0\) such that
\(d_0^{-1}(F)\cap d_1^{-1}(z)\) is nonempty,
\(d_0(d_0^{-1}(F)\cap d_1^{-1}(z))\) is finite, and the restriction of
\(d_1\) to \(d_0^{-1}(F)\) is flat at the points of \(v^{-1}(z)\).
We can therefore resume the notation of
Lemma~\hyperref[lem:sga3-V-7.3]{7.3}, writing \(F_1\) and \(F_2\) for
the inverse images of \(F\) in \(X_1\) and \(X_2\), and so on:
\SGAThreeDiagram[0.56\textwidth]
  {figures/exp5/exp5_localp30_F_groupoid.png}
  {The induced diagram on \(F_2,F_1,F\) used in the proof of \hyperref[thm:sga3-V-8.1]{Theorem~8.1}.}
  {fig:sga3-V-8.1-F-groupoid}
  {LNM 151, pp.~282--283}

As in \hyperref[lem:sga3-V-7.3]{7.3}, one proves that
\(\widetilde d_1\) is quasi-finite at the points of
\(\widetilde d_1^{-1}(z)\). It is therefore natural to consider the
open subset \(F'_1\subseteq F_1\) consisting of the points at which
\(\widetilde d_1\) is both flat and quasi-finite. By
Lemma~\hyperref[lem:sga3-V-7.4]{7.4}, the inverse images of \(F'_1\)
under \(\widetilde d'_1\) and \(\widetilde d'_0\) both consist of the
points of \(F_2\) at which \(\widetilde d'_2\) is flat and
quasi-finite. Thus these inverse images coincide, and \(F'_1\) is of
the form \(\widetilde d_0^{-1}(F')\) for an open subset \(F'\subseteq F\)
(SGA~1, VIII, \SGARefLink{sga3:viii:corollary:4:4}{4.4}). After replacing \(F\) by \(F'\), we may therefore
suppose that \(\widetilde d_1\) is quasi-finite and flat. Let \(W_z\)
be the largest open subset of \(\widetilde d_1(F_1)\) over which
\(\widetilde d_1\) is finite and flat.

The open subset \(W_z\) need not contain \(z\), but it contains the
generic points of the irreducible components of \(X_0\) passing
through \(z\).%
Since \(d_0\) (and likewise \(d_1\)) is faithfully flat and of finite
presentation, hence open, SGA~1, VIII, \SGARefLink{sga3:viii:corollary:5:7}{5.7} shows that
\(d_0^{-1}(W_z)\) and \(d_1^{-1}(W_z)\) both coincide with the largest
open subset of \(\widetilde d'_2(F_2)\) over which
\(\widetilde d'_2\) is finite and flat. As in
\hyperref[lem:sga3-V-7.3]{7.3}, the inverse images of
\(\widetilde d_1^{-1}(W_z)\) under \(\widetilde d'_0\) and
\(\widetilde d'_1\) therefore coincide. Hence
\(\widetilde d_1^{-1}(W_z)=\widetilde d_0^{-1}(U)\) for an open subset
\(U\subseteq F\), and \(U\) is a quasi-section for the groupoid induced
by \(X_\ast\) on \(W_z\).

\section*{9. Elimination of the Noetherian hypotheses in \hyperref[thm:sga3-V-7.1]{Theorem 7.1}}
\phantomsection\label{sec:sga3-V-9}
\addcontentsline{toc}{section}{9. Elimination of the Noetherian hypotheses in \texorpdfstring{\protect\hyperref[thm:sga3-V-7.1]{Theorem 7.1}}{\SGARefLink{thm:sga3-V-7.1}{Theorem 7.1}}}

\medskip
\phantomsection\label{par:sga3-V-9-a}%
\noindent\textbf{a)}\quad
Resume the notation and hypotheses of
Lemma~\hyperref[lem:sga3-V-6.1]{6.1}, and let
\(\pi:S'\to S\) be an arbitrary base change. If \(f:X\to Y\) is a
morphism of \(S\)-schemes, write \(f':X'\to Y'\) for the morphism of
\(S'\)-schemes obtained from \(f\) by extension of the base along
\(\pi\). With this convention, \(p':X'_0\to Y'\) is surjective, as is
the morphism
\[
  X'_1\longrightarrow X'_0\times_{Y'}X'_0
\]
with components \(d'_0\) and \(d'_1\). Thus the set underlying \(Y'\)
is the quotient of the set underlying \(X'_0\) by the equivalence
relation on \(X'_0\) defined by the \(S'\)-groupoid \(X'_\ast\).
Moreover, \(q':U'\to Y'\) is integral and surjective. The topology of
\(Y'\) is consequently the quotient topology of that of \(U'\), and
hence also of that of \(X'_0\); see the proof in
\hyperref[par:sga3-V-6-c]{\S6(c)}.

On the other hand, \(U'\) is plainly a quasi-section of the
\(S'\)-groupoid \(X'_\ast\), so Lemma~\hyperref[lem:sga3-V-6.1]{6.1}
applies. In particular, \(X'_\ast\) has a cokernel \((Y_1,p_1)\), and
the topological space underlying \(Y_1\) is obtained from that
underlying \(X'_0\) by identifying the points equivalent under the
relation defined by \(X'_\ast\). It follows that the canonical
morphism \(Y_1\to Y'\) is a homeomorphism. In fact, it is a universal
homeomorphism. Indeed, if \(S''\) is an \(S'\)-scheme, let \(Y_2\) be
the cokernel of
\((d_0\times_S S'',d_1\times_S S'')\). The preceding argument,
applied to the base changes \(S''\to S'\) and \(S''\to S\), shows that
\[
  Y_2\longrightarrow Y_1\times_{S'}S''
  \qquad\text{and}\qquad
  Y_2\longrightarrow Y\times_SS''\simeq Y'\times_{S'}S''
\]
are homeomorphisms. Hence
\[
  Y_1\times_{S'}S''\longrightarrow Y'\times_{S'}S''
\]
is a homeomorphism as well.

\medskip
\phantomsection\label{par:sga3-V-9-b}%
\noindent\textbf{b)}\quad
The same remarks plainly apply when the groupoid \(X_\ast\)
\emph{locally} possesses quasi-sections (see the proof of
Theorem~\hyperref[thm:sga3-V-7.1]{7.1}).%
\footnote{\label{note:sga3-V-46}Editor's note: What follows has been
expanded in order to bring out \hyperref[thm:sga3-V-9.0]{Theorem~9.0 below}.}
For example, we have \hyperref[thm:sga3-V-9.0]{the following theorem}:
